\section{Scope, controls, and outstanding tasks}\label{sec:outstanding}

The main results now form a quantitative chain. The factorial construction
specifies the complete observable and its self correction. Source
stability bounds the error in replacing it. The mixed-moment theorem
specifies the further joint information needed for a canonical profile.
The Newman-time theorem transports supplied initial information through
collisions. Finally, exact Gaussian integration turns one declared
projection into a finite acquisition problem with separate spatial,
geometric and arithmetic allowances. None of these steps supplies the
missing canonical joint estimate by itself.

The contribution is measured by the displayed error bounds, not by the
number of representations or formal identities. The standard factorial,
contour and Gaussian identities are used as tools; no literature novelty
is claimed for them. Negative results retain their stated classes:
harmonic replacements on the full planar supports, linear Gaussian
convolution passed through a logarithm, and the particular information
lost by a difference-only projection.

The collision controls catalogued in the repository were screened
analytically against these claims; their earlier numerical
certificates were not rerun for this manuscript.
Generic identities and approximation estimates are expected to
hold for controls satisfying their hypotheses. The missing
canonical premise in the following table explains why a control
does not become a theta example; it does not establish a
favorable theta estimate.

\begin{center}\small
\begin{tabularx}{\textwidth}{@{}>{\raggedright\arraybackslash}p{0.29\textwidth}Y@{}}
\toprule
Control family & Premise absent from a canonical inference\\
\midrule
Euler-factor modifications &
The unchanged complete theta coefficients and conductor. Preserved
real zeros do not account for added nonreal zeros.\\
Weighted-square sources and finite coefficient or tail changes &
The complete normalized reflection and unchanged arithmetic
source. Finite positivity or endpoint agreement is insufficient.\\
Reflected super-Gaussians &
The integer-square theta arithmetic, despite positivity and
reflection symmetry.\\
Smooth shape and positive-primitive controls &
The full canonical coefficient expansion, despite matching
shape or auxiliary-transform properties.\\
Finite-prime orbit controls &
The full reflection and all-coefficient source identity;
finite prime data do not determine the residual divisor.\\
Atomic event and finite-profile controls, with continuous extensions &
The full theta profile and coefficient structure; finite
observations do not certify the complete integer zero divisor.\\
RH-conditional normalized endpoint controls &
Ordinary Dirichlet-series arithmetic. The full correlation retains
added double zeros with multiplicity; no simplicity premise is used.\\
Variable observation kernels &
The cosine oscillator and the corresponding Newman equation;
fixed amplitude does not restore the altered observation law.\\
Actual-theta three-profile relaxation &
The coupling of all observations to an integer zero divisor;
the relaxation permits fractional masses.\\
\bottomrule
\end{tabularx}
\end{center}

The polynomial controls in
\eqref{eq:pair-control} and \eqref{eq:real-divisor-control}
are additional local information tests, not members of the
unchanged theta family. The catalog and the project's scope
rules are cited in the accompanying manuscript guide and
continuation note, with source paths rather than new claims
of universal obstruction.

The outstanding tasks are the following, in order of analytical importance.

\begin{enumerate}[leftmargin=*,itemsep=0.55em]
\item \textbf{Acquire genuine same-center canonical information.}
For the fixed compact regime \eqref{eq:fixed-regime}, obtain the moments
in \eqref{eq:joint-moments}, or a sharper estimate of their contracted
combination, with the explicit error in Theorem~\ref{thm:joint-interface}.
For the Gaussian regime, the analogous target is the joint field in
\eqref{eq:gaussian-log-contraction}. A one-point estimate or another
identity for the covariance does not satisfy this task. A paid full
profile is a worthwhile result independently of a sine law or an RH
application.

\item \textbf{Make the time allowance informative for supplied data.}
Theorem~\ref{thm:time-modulus} is complete at fixed compact window and
positive blur. A practical transport application still needs an initial
profile, a quantitative value or usable bound for $A$, and a time interval
on which the additive allowance is useful. Extending this theorem to a
Gaussian physical window requires summing the moving-zero and exterior
costs with its tails. It is not obtained merely by substituting the
Gaussian window into the compact constants. Time-dependent windows or
unfolding scales also require their own derivative terms.

\item \textbf{Extend quantitative transfer to higher orders and limits.}
The hierarchy is defined at every finite order, but higher canonical
error bounds must subtract every repeated-label partition and control
all Gaussian product derivatives. Growing order requires a simultaneous
order-tail bound. Removing blur returns a measure, not a pointwise
smooth density; growing height changes the source and time constants.
Normalization by factorial mass or first intensities requires positive
lower bounds and errors for those denominators.

\item \textbf{Complete the declared finite projection as supporting evidence.}
The inherited census and Theorem~\ref{thm:cell-ranges} pay the geometric
allowance for $P_{1000}$. The endpoint evaluation still has to return
outward enclosures with $r_L+r_U\le1/1000$ and include the complete
$2^{-32}$ tail. The fixed E13 r2 contract permits that bounded evaluation;
this manuscript supplies no scalar value. A failure of that arithmetic
gate would be an execution result, not a failure of the analytical
width theorem or of all centroid methods.

\item \textbf{Supply an additional consumer before drawing collision conclusions.}
A regular full profile can coexist with a collision, as the quadratic
heat control shows. Any collision application needs a resolving
observation scale and a proved lower signal, together with an upper or
signed comparison that beats it. Conditioning, an initial gap, or a
signed exterior estimate must be stated and paid where used. The
collision-constraint investigation owns that separate priority.

\item \textbf{Assess the integrated claims at their actual evidence levels.}
The E07, E08 and E12 acceptances retain their original scopes.
The Gaussian-window and acquisition arguments, the arbitrary-separation
form of \eqref{eq:range-width}, and this assembly have proof-author
checking only. A separate claim audit should check their interfaces
and constants. The earlier B02 and B04 assessments E03 and E04 remain
pending; successful typesetting and finite bookkeeping do not replace them.
\end{enumerate}

The appendices preserve the useful comparison results and their limits,
then record the finite evidence and its outstanding arithmetic gate.
